\phantomsection
\addcontentsline{toc}{section}{ВВЕДЕНИЕ}

\section*{Введение}

Формирование расписания учебных занятий является одной из основных и наиболее сложных задач автоматизации управления учебным процессом вуза. Каждый семестр специальный отдел в составе Учебно-методического управления МГТУ «СТАНКИН» \cite{stankin:schedule} составляет расписание на следующие полгода, выкладывает его на сайт университета и отправляет преподавателям и старостам групп. На сегодняшний день процесс создания расписания учебных занятий производится вручную, что является медленным и ресурсозатратным действием.

\textbf{Целью} выпускной квалификационной работы является повышение эффективности деятельности учебного заведения за счёт исследования методов автоматической генерации расписания и разработки программного средства, реализующего такую генерацию с учётом регламента вуза, предпочтений преподавателей и «компоновки вуза».

\textbf{Объектом} исследования является программная поддержка процесса создания расписания учебных занятий.

\textbf{Предметом} исследования являются архитектура и функционирование клиент-серверного веб-приложения для создания расписания учебных занятий с учётом личных предпочтений преподавательского состава.

Для достижения цели поставлены следующие \textbf{задачи}:
\begin{enumerate}
	\item провести исследование методов автоматического составления расписаний учебных занятий, обосновать необходимость собственного подхода и проанализировать правила составления расписания МГТУ «СТАНКИН»;
	\item сформулировать требования к клиент-серверному веб-приложению для создания расписания учебных занятий;
	\item разработать проект клиент-серверного веб-приложения;
	\item реализовать клиент-серверное веб-приложение в соответствии с разработанным проектом.
\end{enumerate}

\textbf{Методы исследования:} системный анализ, процессно-ориентированный подход и функциональное моделирование процессов, объектно-ориентированный анализ, сравнительный анализ алгоритмов и программная реализация предложенного решения.

В работе будут рассмотрены методы автоматического создания расписания, предложен гибридный алгоритм на основе коллапса волновой функции (WFC) и эволюционной оптимизации, а также реализован сервис, выполняющий следующие функции:
\begin{enumerate}
	\item генерация расписания с учётом учебных планов, нормативов, списков преподавателей, групп и т.~д.;
	\item настройка распределения занятий по неделе и в течение дня, в том числе формирование разгрузочных дней и смещение учебного процесса в пределах суток;
	\item «компоновка вуза» для сокращения времени на переходы между учебными занятиями;
	\item вывод расписания в различных форматах, включая веб-интерфейс для просмотра преподавателями и студентами.
\end{enumerate}

Так как качество расписания определяет эффективность учебного процесса, значительное внимание уделено оценке расписаний и их последующему улучшению.

\textbf{Научная новизна} заключается в адаптации WFC к модели «группа~$\times$~день~$\times$~слот» с учётом регламента вуза и в гибридной схеме оптимизации: построение допустимого по жёстким ограничениям расписания с последующим улучшением эволюционным методом с repair-мутациями и отдельным этапом назначения аудиторий.

\textbf{Практическая ценность} состоит в автоматизации ресурсозатратного процесса планирования учебных занятий на семестр и в снижении числа нарушений нормативов при пересборке расписания.

\textbf{Структура работы.} Первая глава содержит обзор методов решения задачи. Во второй главе сформулированы требования к расписанию и описан предложенный алгоритм. Третья и четвёртая главы посвящены архитектуре и выбору средств реализации. Пятая глава описывает программную реализацию и результаты экспериментов. Экономическое обоснование приведено в отдельной главе. В заключении подведены итоги и намечены направления дальнейшей работы.
